\documentclass[11pt,letterpaper]{article}
\usepackage[T1]{fontenc}
\usepackage{lmodern}
\usepackage[margin=1in]{geometry}
\usepackage{amsmath,amssymb,amsthm,mathtools,microtype}
\usepackage{xcolor,hyperref}
\hypersetup{colorlinks=true,linkcolor=blue!45!black,citecolor=blue!45!black,
urlcolor=blue!45!black,pdftitle={Vacuum-relative stability and charge gaps in compact Hamiltonian lattice QED},
pdfauthor={}}
\setlength{\emergencystretch}{2em}
\numberwithin{equation}{section}
\newtheorem{theorem}{Theorem}[section]
\newtheorem{lemma}[theorem]{Lemma}
\newtheorem{proposition}[theorem]{Proposition}
\newtheorem{corollary}[theorem]{Corollary}
\theoremstyle{remark}
\newtheorem{remark}[theorem]{Remark}
\DeclareMathOperator{\Tr}{Tr}
\DeclareMathOperator{\diver}{div}
\newcommand{\F}{\mathcal F}
\newcommand{\G}{\mathcal G}
\newcommand{\vac}{\mathrm{vac}}
\title{Vacuum-relative stability and charge gaps\\in compact Hamiltonian lattice QED}
\author{}
\date{29 October 2026}
\begin{document}
\maketitle
\begin{center}
\small Research-tool credits: Opus 5.5 + Astra 6\\
Human author metadata has not been supplied.
\end{center}
\begin{abstract}
For a precisely specified compact $U(1)$ Hamiltonian lattice gauge theory
with four-component Wilson fermions, nonnegative bare mass $m$, and
untruncated electric rotors, we prove
$E_\nu\ge E_0+m|\sum_x\nu_x|$ in every admissible integer static-source
sector of exact local Gauss law. The comparison holds at every positive
gauge coupling. It includes the sources' electric self-energy;
subtracting prescribed individual nuclear self-energies gives a
volume-uniform lower bound of order the number of nuclei.
The argument combines a positive centered fermion trace, a quantitative
transfer-spectrum estimate, positive compact-rotor heat kernels, and
gauge projection. All transfer and Fourier estimates are proved here.
At fixed lattice spacing we also prove the vacuum thermodynamic limit
and an open, zero-exterior-flux thermodynamic limit for each fixed
admissible periodic source cell. The corresponding periodic-boundary
source limit requires a boundary-flux comparison that is stated
explicitly and remains unproved. These are finite-spacing results;
they do not construct renormalized continuum QED or establish a
neutral pair gap, binding, or crystalline order.
\end{abstract}

\paragraph{Tool-use disclosure.}
Significant portions of the research development, proof drafts,
exposition, and internal checking were produced with generative-AI
tools. Internal AI-assisted audits are not external peer review.

\section{Introduction}
An absolute lower bound proportional to the number of lattice sites
does not compare matter with the interacting vacuum. The question
addressed here is whether introducing static integer charges can
lower the energy below that vacuum when all fermion pair terms and
the quantized gauge field are retained. For the Hamiltonian below,
the answer is negative before the individual source self-energies
are subtracted. The induced net fermion charge also costs at least
its bare rest energy.

The compact-link framework is that of Hamiltonian lattice gauge
theory \cite{KS}. We use a direct spatial Wilson discretization of
the Dirac Hamiltonian \cite{Zache}, with continuous time and canonical
fermions. Its positive Wilson mass coefficient is a structural
hypothesis of the proof. Positivity criteria for fermion traces
\cite{Wei} and Gaussian-domination estimates in electron--radiation
models \cite{Miyao} provide related context; no theorem for those
different Hamiltonians is substituted for the proof here.
Likewise, continuum Coulomb thermodynamic-limit results
\cite{ThermoI,ThermoII} do not settle the compact Gauss-law boundary
problem encountered below. No priority claim is made.

The source character in our gauge projection has modulus one, so
its integral against a positive measure is bounded by the same
integral without that character. Such a modulus observation does
not by itself identify a static energy or fix its self-energy
normalization. Seiler's lattice result \cite{Seiler} concerns an
upper bound on growth of the color-confining potential using
Hamiltonian positivity and Euclidean rotations; it is not invoked
as the source-sector theorem proved here.
The Gaussian-domination comparison of
Fr\"ohlich--Simon--Spencer \cite[Section 2]{FSS} concerns shifted
bond energies in classical ferromagnets. These are relevant
predecessors in method, with different hypotheses and conclusions.
Here the finite-CAR positivity, gauge-sector projection and
zero-temperature passage are all established explicitly.

Section~\ref{sec:model} specifies the operator, ordering, gauge
conventions and admissible sectors. Section~\ref{sec:transfer}
proves the finite-dimensional positivity and Fourier estimates.
Section~\ref{sec:comparison} applies them to exact gauge projection
and fixes the nuclear self-energy convention.
Section~\ref{sec:thermo} proves the two thermodynamic limits and
separates the remaining periodic-source question.

\section{Hamiltonian and exact Gauss sectors}
\label{sec:model}

Let $X=(\mathbb Z/n\mathbb Z)^3$, $n\ge3$, have spacing $a>0$ and
$N=n^3$ sites. Each positively oriented link $e=(x,j)$ joins $x$
to $x+e_j$; write $\mathcal E$ for the $3N$ links. The auxiliary
Hilbert space is
\begin{equation}
 \mathcal H=L^2(U(1)^{\mathcal E},dU)
       \otimes\F_a(\ell^2(X)\otimes\mathbb C^4).
 \label{eq:hilbert}
\end{equation}
All Haar measures are normalized. On a link, $U_e=e^{i\varphi_e}$
is multiplication and $E_e=-i\partial_{\varphi_e}$ has spectrum
$\mathbb Z$. No electric-field truncation is imposed.
For a coupling $g>0$, put
\begin{equation}
 H_E=\frac{g^2}{2a}\sum_eE_e^2,\qquad
 H_B(U)=\frac1{g^2a}\sum_p(1-\operatorname{Re}U_p).
 \label{eq:gauge}
\end{equation}
Each elementary plaquette is counted once, with
$U_p=U_{xj}U_{x+e_j,k}\overline U_{x+e_k,j}\overline U_{xk}$,
$j<k$. There is no background electric shift, topological angle,
or additional flux-sector projection.

Let $\alpha_1,\alpha_2,\alpha_3,\beta$ be Hermitian $4\times4$
matrices that square to one and anticommute pairwise. For negative
fermionic electric charge use the link convention
\begin{align}
 (T_j(U)f)(x)&=\overline U_{xj}f(x+e_j),&
 v_j(U)&=\frac{T_j(U)-T_j(U)^*}{2ia},\nonumber\\
 M(U)&=m+\frac r{2a}\sum_{j=1}^3(2-T_j(U)-T_j(U)^*),&
 h_W(U)&=\sum_j\alpha_jv_j(U)+\beta M(U),
 \label{eq:wilson}\\
 H&=H_E+H_B(U)+d\Gamma(h_W(U)).&&
 \label{eq:hamiltonian}
\end{align}
The bare parameters satisfy $m\ge0$ and $r\ge0$. Spatial matrices
commute with the Dirac matrices. Since $T_j(U)$ is unitary,
\begin{equation}
 M(U)\ge mI,\qquad
 \|h_W(U)\|\le h_*:=m+(3+6r)/a.
 \label{eq:mass-bound}
\end{equation}
The first assertion follows from $2-T-T^*=(I-T)^*(I-T)$.
For $r>0$ this is the usual positive Wilson term; $r=0$ also obeys
the proof but retains the naive doublers.

The fermion term in \eqref{eq:hamiltonian} is the \emph{raw bilinear}.
There is no link-dependent vacuum subtraction. In a particle--hole
representation relative to the negative spectral subspace of
$h_W(1)$, the scalar $\Tr(P_-(h_W(1))h_W(U))$ is retained.
Subtracting that function of $U$ changes the model. A common
source-independent scalar constant may be changed without affecting
any relative conclusion.
For $m>0$, $h_W(1)$ has $2N$ positive and $2N$ negative levels.
The particle--hole transform identifies electron excess with
$N_e-N_p$; the blocks $P_+h_W(U)P_-$ give electron--positron
creation and annihilation. They are not removed.

Write $n_x=\sum_{s=1}^4c_{xs}^*c_{xs}$ and $q_x=n_x-2$.
The physical fermionic charge at $x$ is $-gq_x$.
For a gauge function $u_x=e^{i\theta_x}$ define
\begin{equation}
 U^u_{xj}=u_xU_{xj}\overline u_{x+e_j},\qquad
 (R_u\Psi)(U)=\Psi(U^u),\qquad
 W(u)=e^{i\sum_x\theta_xq_x}R_u.
 \label{eq:gauge-unitary}
\end{equation}
Then $h_W(U^u)=u^*h_W(U)u$, so $W(u)$ commutes with each term of
$H$. Its strongly commuting integer generators are
\begin{equation}
 G_x=\sum_j(E_{xj}-E_{x-e_j,j})+q_x.
 \label{eq:gauss}
\end{equation}
For an integer static source $\nu:X\to\mathbb Z$, let
\begin{equation}
 \mathcal H_\nu=\{\Psi:G_x\Psi=\nu_x\Psi\ \forall x\},\qquad
 P_\nu=\int\frac{d\theta}{(2\pi)^N}
          e^{-i\theta\cdot\nu}W(e^{i\theta}).
 \label{eq:projection}
\end{equation}
This is an orthogonal projection. Gauss's law is
$\diver E=\nu-q$, so positive sources require electron excess.
Writing $Q_\nu=\sum_x\nu_x$, one has exactly
\begin{equation}
 \sum_xq_x=Q_\nu,\qquad
 \mathcal H_\nu\ne0\ \Longleftrightarrow\ |Q_\nu|\le2N.
 \label{eq:capacity}
\end{equation}
Necessity follows by summing Gauss's law. Conversely choose
occupations with that total excess. The integer function $\nu-q$
has zero sum and is the divergence of an integer flow on a
spanning tree. Its electric Fourier vector and the occupation
vector give a physical state. This also shows why integrality
and finite fermion capacity cannot be suppressed.

\begin{lemma}[Operator realization]
\label{lem:operator}
$H$ is self-adjoint on $D(H_E)$, has compact resolvent, and
$e^{-bH}$ is trace class for every $b>0$. Every nonempty
$\mathcal H_\nu$ reduces $H$ and has these properties. Moreover
\begin{equation}
 H\ge-2Nh_*,\qquad
 E_0\le-2N(m+3r/a)+3N/(g^2a),
 \label{eq:extensive}
\end{equation}
where $E_\nu=\inf\sigma(H|_{\mathcal H_\nu})$.
\end{lemma}
\begin{proof}
$H_B+d\Gamma(h_W)$ is a bounded continuous Hermitian matrix
multiplier on the compact link torus. Bounded perturbation of its
positive Laplacian proves the operator assertions; eigenvalue
comparison proves trace-class heat semigroups.
The Hermitian unitary $C_5=\alpha_1\alpha_2\alpha_3\beta$
anticommutes with $h_W(U)$. Its spectrum is paired under
$\lambda\mapsto-\lambda$, so
$d\Gamma(h_W(U))\ge-\frac12\Tr|h_W(U)|\ge-2Nh_*$.
For the upper bound use the constant link wave function, with the
two negative-$\beta$ spinors occupied at every site. It has
$q_x=E_e=0$, hence is in $\mathcal H_0$.
Hopping expectations vanish; its fermion energy is
$-2N(m+3r/a)$ and magnetic energy is at most $3N/(g^2a)$.
\end{proof}

The physical space is obtained by projection from \eqref{eq:hilbert}.
No factorization into independent physical matter, longitudinal,
and transverse Hilbert spaces is assumed.

\section{Positive traces and a quantitative charge estimate}
\label{sec:transfer}

The finite-dimensional argument allows spatial coefficients and
the mass matrix to vary in time, as required for quantized link
histories. In this section the one-particle dimension is $D=2d$,
with the same four-component Dirac structure as above.

\begin{lemma}[Positivity and the transfer annulus]
\label{lem:transfer}
Consider ordered free evolutions generated by
\[
 h(t)=\sum_j\alpha_jv_j(t)+\beta M(t),\qquad M(t)\ge m_*I>0,
\]
with bounded Hermitian spatial coefficients and total free time
$b>0$, interspersed with spin-scalar unitary kicks.
Let $U$ be their one-particle product. Include each kick's
centering phase in the finite-CAR trace
\begin{equation}
 F=e^{\frac12\int\Tr A(t)\,dt}\det(I+U),
 \qquad U'=-A(t)U,\quad A=c(t)h(t)+i\Phi(t).
 \label{eq:centered-trace}
\end{equation}
The auxiliary ODE interval is $[0,T]$; free and kick intervals
can be separate, $c\ge0$ obeys $\int_0^Tc(t)\,dt=b$, and
$\Phi$ is Hermitian and spin scalar.
Then $F>0$. Exactly $d$ transfer eigenvalues, with algebraic
multiplicity, obey
\begin{equation}
 |\lambda|\le e^{-bm_*};
 \label{eq:annulus}
\end{equation}
the remaining $d$ are their partners $1/\overline\lambda$ and
have modulus at least $e^{bm_*}$. Normality or diagonalizability
of $U$ is not required.
\end{lemma}
\begin{proof}
Since $C_5$ anticommutes with $h$ and commutes with $\Phi$,
\[
 A^*C_5+C_5A=0,\qquad U^*C_5U=C_5.
\]
Thus $U^*=C_5U^{-1}C_5$ and the spectrum is paired, with
algebraic multiplicities, under $\lambda\mapsto1/\overline\lambda$.
Also
\[
 \overline{\det(I+U)}=\det(U)^{-1}\det(I+U),\qquad
 \det U=e^{-\int\Tr A}.
\]
Because $\Tr h=0$, $\int\Tr A$ is purely imaginary. These
identities make \eqref{eq:centered-trace} real.

For any transfer eigenvector $v$, transport it along the
successive intervals and kicks, obtaining $f(t)$.
The form $B(t)=\langle f(t),\beta f(t)\rangle$ is unchanged by
scalar kicks and satisfies
\begin{equation}
 B'(t)=-2c(t)\langle f(t),M(t)f(t)\rangle
       \le-2m_*c(t)\|f(t)\|^2.
 \label{eq:beta-form}
\end{equation}
Its endpoint is $B(T)=|\lambda|^2B(0)$, and its total decrease
is strict because $b>0$ and a nonzero solution never vanishes.
Hence $|\lambda|=1$ is impossible.
If $|\lambda|<1$, both endpoints and the entire trajectory
have $B>0$. Since $\|f\|^2\ge B$, integration of
$B'\le-2m_*c(t)B$ proves \eqref{eq:annulus}.
If $|\lambda|>1$, the entire trajectory has $B<0$;
$\|f\|^2\ge-B$ gives $B'\le2m_*c(t)B$ and
$|\lambda|\ge e^{bm_*}$, reversing the inequality when
dividing by $B(0)<0$. Spectral pairing gives exactly $d$
eigenvalues in each region.

In particular $\det(I+U)\ne0$. Continuously turn off all kicks
and all $v_j$, and replace $M(t)$ by
$(1-s)M(t)+sm_*I$. The real centered trace never vanishes
along this deformation. At its endpoint it equals
$\det(I+e^{-bm_*\beta})>0$, which proves its sign.
The same argument holds for bounded measurable histories:
their solutions are absolutely continuous and
\eqref{eq:beta-form} holds almost everywhere.
\end{proof}

For a scalar kick $u=e^{i\Theta}$ the centering factor is
$e^{-\frac i2\Tr\Theta}$. For a local gauge twist,
$\Tr\Theta=4\sum_x\theta_x$, so it is precisely
$e^{-2i\sum_x\theta_x}$ from $q_x=n_x-2$.
This phase is indispensable and is single-valued when
$\theta_x$ changes by $2\pi$.
For a nonnegative mass coefficient, replace $M$ by
$M+\epsilon I$ and pass to the finite-dimensional limit.
The centered trace is then nonnegative but may vanish.

\begin{lemma}[Canonical Fourier bound]
\label{lem:fourier}
Append to a history of Lemma~\ref{lem:transfer} a global
charge twist $e^{i\phi Q}$, $Q=d\Gamma(I)-d$.
Let $F(\phi)$ be its centered trace and
$F_q=(2\pi)^{-1}\int_0^{2\pi}e^{-iq\phi}F(\phi)\,d\phi$.
For $|q|\le d$,
\begin{equation}
 |F_q|\le\binom{2d}{d+|q|}e^{-bm_*|q|}F_0.
 \label{eq:fourier-bound}
\end{equation}
\end{lemma}
\begin{proof}
Let $\lambda_1,\ldots,\lambda_d$ be the transfer eigenvalues
in the small disk. Pairing determinant factors and including
the phase $e^{-id\phi}$ gives
\[
 F(\phi)=C\prod_{j=1}^d|1+e^{i\phi}\lambda_j|^2,\qquad C>0.
\]
A paired factor is
$e^{i\phi}\overline\lambda_j^{-1}
 |1+e^{i\phi}\lambda_j|^2$; the remaining constant is positive
because $F(0)>0$. This uses determinant eigenvalues only and
also holds for defective matrices.
If $e_j$ are the elementary symmetric polynomials in the
$\lambda_j$ and $s=e^{-bm_*}$, then
\[
 F_0=C\sum_{j=0}^d|e_j|^2\ge C,\qquad
 |e_j|\le\binom dj s^j.
\]
For $q\ge0$, Vandermonde's convolution yields
\[
 |F_q|\le Cs^q\sum_{j=0}^{d-q}\binom d{j+q}\binom dj s^{2j}
 \le Cs^q\binom{2d}{d+q}
 \le s^q\binom{2d}{d+q}F_0.
\]
Negative $q$ follows by complex conjugation.
\end{proof}

The prefactor in \eqref{eq:fourier-bound} need not be one.
It is retained until division by inverse temperature in the
ground-energy limit.

\section{Vacuum comparison and nuclear self-energies}
\label{sec:comparison}

\begin{theorem}[Exact physical-sector comparison]
\label{thm:main}
For the Hamiltonian \eqref{eq:hamiltonian}, with $a>0$, $g>0$,
$r\ge0$, and $m\ge0$, every admissible integer source satisfies
\begin{equation}
 E_\nu\ge E_0+m|Q_\nu|,\qquad Q_\nu=\sum_x\nu_x.
 \label{eq:main}
\end{equation}
At inverse temperature $b>0$ let $Z_\nu(b)=\Tr(P_\nu e^{-bH})$.
Then $0<Z_\nu(b)\le Z_0(b)$, and for $m>0$ also
\begin{equation}
 Z_\nu(b)\le\binom{4N}{2N+|Q_\nu|}
                         e^{-bm|Q_\nu|}Z_0(b).
 \label{eq:thermal}
\end{equation}
The unrestricted auxiliary ground energy equals the physical
zero-source energy $E_0$.
\end{theorem}
\begin{proof}
Set $c_E=g^2/(2a)$. The electric heat kernel on one link,
relative to normalized Haar measure, is the positive wrapped Gaussian
\begin{equation}
 k_t(\phi)=\sum_{\ell\in\mathbb Z}e^{-c_Et\ell^2}e^{i\ell\phi}
 =\sqrt{\frac{\pi}{c_Et}}\sum_{s\in\mathbb Z}
             e^{-(\phi-2\pi s)^2/(4c_Et)}>0.
 \label{eq:heat}
\end{equation}
Let $K_t(U,V)$ be the product over links. The bounded matrix
Feynman--Kac formula in this setting follows directly from the
Dyson series. With $B(U)=H_B(U)+d\Gamma(h_W(U))$, the norm of
the $k$th-order kernel is bounded by
\[
 K_b(U,V)(b\|B\|)^k/k!.
\]
To see this, bound every matrix insertion by $\|B\|$ and
integrate the positive kernels by their convolution identity.
The series is absolutely convergent and is the bounded-perturbation
Dyson series for the semigroup. Its Markov bridge moments give
\begin{equation}
 [e^{-bH}](U,V)=K_b(U,V)\,
 \mathbb E_{V\to U}^{b}
 \left[e^{-\int_0^bH_B(\omega_t)\,dt}\Gamma(V_\omega)\right],
 \label{eq:fk}
\end{equation}
where $V_\omega$ is the ordered one-particle propagator generated
by $-h_W(\omega_t)$. The circle Brownian bridges have continuous
paths. Boundedness of the coefficients justifies step approximations
and dominated convergence. This establishes the formula without
an unbounded-potential or formal Trotter interchange.

Inserting the gauge unitary changes the first kernel argument
to $U^u$. The finite-CAR identity $\Tr_\F\Gamma(A)=\det(I+A)$ gives
\begin{align}
 \Tr(W(u)e^{-bH})
 =\int dU\,K_b(U^u,U)\,
 \mathbb E_{U\to U^u}^{b}
 \left[e^{-\int_0^bH_B(\omega_t)\,dt}
 e^{-2i\sum_x\theta_x}\det(I+uV_\omega)\right].
 \label{eq:twisted}
\end{align}
The trace formula is justified by factoring the heat semigroup
into two Hilbert--Schmidt factors and integrating their kernel
convolution. The preceding Dyson bound also dominates the
integrals uniformly in the gauge angles at fixed finite lattice
and fixed $b$.

For $m>0$ every integrand in \eqref{eq:twisted} is positive:
the scalar factors are positive and Lemma~\ref{lem:transfer}
applies to the time-dependent Wilson mass $M(\omega_t)\ge mI$.
Denote the resulting nonnegative function of the angles by
$F_b(\theta)$. Exact gauge projection gives
\[
 Z_\nu(b)=\int\frac{d\theta}{(2\pi)^N}
          e^{-i\theta\cdot\nu}F_b(\theta),\qquad
 |Z_\nu(b)|\le Z_0(b).
\]
The actual heat trace of a nonempty sector is positive.

For the stronger assertion, fix a root site and write
$\theta_x=\vartheta_x+\phi$, with $\vartheta_{\rm root}=0$.
This is a measure-preserving change of torus coordinates.
The endpoint $U^u$ and the entire bridge weight depend only
on $\vartheta$. At fixed relative twist and bridge, integration
in $\phi$ is the $Q_\nu$ Fourier coefficient of the centered
trace for $e^{i\vartheta}V_\omega$. Lemma~\ref{lem:fourier}
applies with $D=4N$, $d=2N$, and $m_*=m$.
Remove the remaining phase $e^{-i\vartheta\cdot\nu}$ by absolute
value and integrate the nonnegative zeroth coefficient.
Its integral is exactly the vacuum's full local-Gauss partition
function. This proves \eqref{eq:thermal}.

Compact resolvent gives
$E_\nu=-\lim_{b\to\infty}b^{-1}\log Z_\nu(b)$;
the binomial factor disappears in this limit.
For $m=0$, determinants are nonnegative by the limit in
Section~\ref{sec:transfer}; alternatively
$\|H_\epsilon-H_0\|\le2N\epsilon$ gives continuity of every
sector ground energy. Finally the auxiliary space is the direct
sum of joint integer Gauss sectors. Their comparison with
$\mathcal H_0$ proves its stated global minimality.
\end{proof}

The inequality concerns the net fermion charge forced by the
external sources. It is not a gap above the vacuum within its
neutral sector and says nothing about a neutral pair or photon
excitation gap. The constant gauge transformation enforces
$\sum q_x=Q_\nu$; every relative gauge angle is also projected.
All compact link holonomies, including torus harmonic variables,
remain integrated.

\subsection{The subtraction convention}

Let $D:\mathbb R^{\mathcal E}\to\mathbb R^X$ be the incidence
matrix, $\Delta=DD^*$ the nonnegative combinatorial Laplacian,
and $\G_n=\Delta^+$ its zero-mean inverse. On $\mathcal H_\nu$,
\begin{equation}
 H_E=c_E\left[(\nu-q)^*\G_n(\nu-q)
              +\|E-D^*\G_n(\nu-q)\|^2\right].
 \label{eq:longitudinal}
\end{equation}
This follows by orthogonal projection of a real link vector onto
$\operatorname{ran}D^*$ and $\ker D$, using $DE=\nu-q$.
The electric and occupation operators commute, so the same
calculation holds in their joint basis as a quadratic-form
identity. The allowed spectrum of the residual flux depends
on the charge and integrality constraints; \eqref{eq:longitudinal}
does not create an independent transverse physical Hilbert space.

For nuclei at sites $R_i$, with integer $1\le Z_i\le Z_{\max}$,
put $\nu_x=\sum_iZ_i\mathbf1_{\{x=R_i\}}$.
The source-source term in \eqref{eq:longitudinal} is
\begin{equation}
 c_E\nu^*\G_n\nu
 =c_E\G_n(0)\sum_iZ_i^2
   +2c_E\sum_{i<j}Z_iZ_j\G_n(R_i-R_j).
 \label{eq:source-energy}
\end{equation}
Thus define the prescribed individual subtraction and energy by
\[
 S_n=c_E\G_n(0)\sum_iZ_i^2,\qquad
 E_\nu^{\rm sub}=E_\nu-S_n.
\]
It is state independent and position independent, with its
finite-box dependence fixed by the periodic Green function.
There was no subtraction in Theorem~\ref{thm:main}.

\begin{lemma}[Uniform lattice self-energy]
\label{lem:green}
For every $n\ge3$,
\begin{equation}
 0\le\G_n(0)=\frac1{n^3}\sum_{k\ne0}
 \frac1{4\sum_{j=1}^3\sin^2(\pi k_j/n)}
 \le\frac{13}{16}.
 \label{eq:green}
\end{equation}
\end{lemma}
\begin{proof}
Choose one reciprocal representative per residue with
$|k_j|\le n/2$. The denominator is at least $16|k|^2/n^2$,
since $\sin(\pi|k_j|/n)\ge2|k_j|/n$.
The shell $\max_j|k_j|=s$ has at most $24s^2+2\le26s^2$
points. Therefore
$\G_n(0)\le26\lfloor n/2\rfloor/(16n)\le13/16$.
\end{proof}

\begin{corollary}[Stability after individual nuclear subtraction]
\label{cor:stability}
For any nuclear sites, including repeated sites, with
$Z_{\rm tot}=\sum_iZ_i\le2n^3$,
\begin{align}
 E_\nu^{\rm sub}-E_0
 &\ge mZ_{\rm tot}-c_E\G_n(0)\sum_iZ_i^2\nonumber\\
 &\ge mZ_{\rm tot}-\frac{13\pi\alpha}{8a}Z_{\max}^2K,
 \qquad g^2=4\pi\alpha.
 \label{eq:stability}
\end{align}
\end{corollary}
\begin{proof}
Subtract $S_n$ from \eqref{eq:main} and use Lemma~\ref{lem:green}.
\end{proof}

The constant is independent of volume, source positions, fermion
pair content and gauge excitations. It has no coupling-smallness
condition. Adding $S_n$ to the raw energy would count the source
self-energy twice.
If two labeled nuclei occupy the same site, their mutual
$2c_EZ_iZ_j\G_n(0)$ term is retained. Subtracting the square
of the combined site charge instead is a different model and
is not an $O(K)$ subtraction for unrestricted site occupancy.
Lattice sources resolve distance only to order $a$; this model
does not preserve the continuum point-source collision singularity.

One can alternatively prescribe volume-independent species constants
using
\begin{equation}
 \G_\infty(0)=\int_{[-1/2,1/2]^3}
 \frac{dt}{4\sum_j\sin^2(\pi t_j)}
 \le\frac{\pi\sqrt3}{8}.
 \label{eq:green-infinite}
\end{equation}
The bound follows by comparison with
$\frac1{16}|t|^{-2}$ and the containing ball of radius $\sqrt3/2$.
Then $S_\infty=c_E\G_\infty(0)\sum_iZ_i^2$ is a different
explicit normalization, and \eqref{eq:main} gives the
corresponding bound with $S_\infty$.

\section{Thermodynamic limits at fixed spacing}
\label{sec:thermo}

All parameters are fixed in this section. We distinguish
unrestricted open-graph minima, physical zero-exterior-flux
open minima, and periodic physical minima.

\subsection{Open graphs and a bounded cutting estimate}

For a rectangular $Y\subset\mathbb Z^3$, keep internal links and
plaquettes. The open fermion matrix is $m\beta+\sum_{e\subset Y}h_e$,
where for $e=(x,y=x+e_j)$
\begin{equation}
 h_e=\frac{\alpha_j}{2ia}
 (\overline U_e|x\rangle\langle y|-U_e|y\rangle\langle x|)
 +\frac{r\beta}{2a}
 (|x\rangle-U_e|y\rangle)(\langle x|-\overline U_e\langle y|).
 \label{eq:open-edge}
\end{equation}
The Wilson diagonal at $x$ is $m+r\deg_Y(x)/(2a)$,
not the periodic value at boundary sites. Its mass coefficient
is still at least $m$. Denote the auxiliary minimum by
$e_{\rm op}(Y)$ and the physical minimum with
$\diver_YE+q=\eta$ by $E_Y^{\rm op}(\eta)$.
There are no exterior link rotors: this convention fixes exterior
flux to zero.

The proof of Theorem~\ref{thm:main} applies on this finite graph:
its incidence matrix replaces the periodic one, the Wilson mass
is positive, and the internal rotor kernels are unchanged.
For connected $Y$ the admissibility condition is
$|\sum_Y\eta|\le2|Y|$, and
\begin{equation}
 e_{\rm op}(Y)=E_Y^{\rm op}(0),\qquad
 E_Y^{\rm op}(\eta)\ge e_{\rm op}(Y)+m|\sum_Y\eta|.
 \label{eq:open-comparison}
\end{equation}
The empty-link case is the same finite-CAR argument.
Uniformly for rectangles of degree at most six,
\begin{equation}
 -2h_*|Y|\le e_{\rm op}(Y)\le3|Y|/(g^2a).
 \label{eq:open-extensive}
\end{equation}
The upper trial has zero electric flux and two negative-$\beta$
spinors at every site; discard its nonpositive fermion energy.

\begin{lemma}[Cutting with retained rotors]
\label{lem:cut}
Select rectangular open blocks partitioning the sites of a cubic
graph. Let $\mathcal C$ be all original links not internal to
these open blocks, including periodic closing links even when
both endpoints lie in the same block, and put $n_c=|\mathcal C|$.
On the original auxiliary space define
\[
 H_{\rm cut}=\sum_B H_B^{\rm op}
                       +c_E\sum_{e\in\mathcal C}E_e^2,
 \qquad c_E=g^2/(2a).
\]
Then
\begin{equation}
 \|H-H_{\rm cut}\|\le Bn_c,\qquad
 B=(4r+2+8/g^2)/a,
 \label{eq:cut-bound}
\end{equation}
and the cut auxiliary minimum equals $\sum_B e_{\rm op}(B)$.
The same statement holds when adjoining periodic closing links
to an open cube, after adding their free rotor factors.
\end{lemma}
\begin{proof}
The endpoint matrix \eqref{eq:open-edge} has norm at most
$r/a+1/(2a)$ and anticommutes with $C_5$. Its eight-dimensional
spectrum is paired, so
$\|d\Gamma(h_e)\|=\frac12\Tr|h_e|\le(4r+2)/a$.
A plaquette multiplier has norm at most $2/(g^2a)$ and a link
meets at most four plaquettes. Remove the cut fermion terms and
every plaquette meeting a cut link; counting a plaquette more
than once only enlarges the estimate.
The unbounded electric term on each cut link is retained as an
independent $L^2(U(1))$ rotor with minimum zero.
Even fermion operators on disjoint blocks commute under the graded
Fock decomposition, so their unrestricted minima add.
The variational principle gives the claimed spectral comparison.
\end{proof}

\subsection{The vacuum limit}

Write $Q_n=\{0,\ldots,n-1\}^3$, $e_{\rm op}(n)=e_{\rm op}(Q_n)$,
and let $E_0^{\rm per}(n)$ be the physical periodic vacuum energy.

\begin{theorem}[Vacuum thermodynamic limit]
\label{thm:vacuum-limit}
The finite limit
\begin{equation}
 e_{\vac}=\lim_{n\to\infty}\frac{e_{\rm op}(n)}{n^3}
         =\lim_{n\to\infty}\frac{E_0^{\rm per}(n)}{n^3}
 \label{eq:vacuum-limit}
\end{equation}
exists, and for every integer $\ell\ge3$,
\begin{equation}
 \left|e_{\vac}-\frac{e_{\rm op}(\ell)}{\ell^3}\right|
 \le\frac{3B}{\ell},\qquad
 \left|e_{\vac}-\frac{E_0^{\rm per}(\ell)}{\ell^3}\right|
 \le\frac{6B}{\ell}.
 \label{eq:vacuum-rate}
\end{equation}
Physical energy per spatial volume is $e_{\vac}/a^3$.
\end{theorem}
\begin{proof}
For $n=k\ell$, cutting an open cube into $k^3$ cubes removes
$3(k-1)n^2\le3n^3/\ell$ links. Hence
$|e_{\rm op}(n)-k^3e_{\rm op}(\ell)|\le3Bn^3/\ell$.
For arbitrary $n$, subdivide each coordinate into intervals of
length $\ell$ and one possible remainder. Remainder boxes contain
at most $3\ell n^2$ sites; their energies obey
\eqref{eq:open-extensive}. The cut-link count is at most
$3n^3/\ell$. Thus at every fixed $\ell$,
\[
 \frac{e_{\rm op}(\ell)}{\ell^3}-\frac{3B}{\ell}
 \le\liminf_n\frac{e_{\rm op}(n)}{n^3}
 \le\limsup_n\frac{e_{\rm op}(n)}{n^3}
 \le\frac{e_{\rm op}(\ell)}{\ell^3}+\frac{3B}{\ell}.
\]
The limsup--liminf difference is at most $6B/\ell$ for every
$\ell$, proving existence and the first estimate.
Closing an open $n$-cube adds $3n^2$ links. Lemma~\ref{lem:cut}
gives $|E_{\rm aux}^{\rm per}(n)-e_{\rm op}(n)|\le3Bn^2$.
Theorem~\ref{thm:main} identifies
$E_{\rm aux}^{\rm per}(n)=E_0^{\rm per}(n)$, completing the proof.
\end{proof}

This cutting proof concerns the auxiliary tensor product.
It does not set the boundary flux of an arbitrary physical
state to zero.

\subsection{A fixed periodic source cell}

Choose a cubic period cell of side $p\ge1$, with $V_c=p^3$ sites,
$K_c\ge1$ labeled positive nuclei and charge $Z_c\le2V_c$.
A rectangular period can be replaced by a cubic supercell.
All periodic source sizes satisfy $n\in p\mathbb N$ and $n\ge3$.
Choose integers $q_x^{(c)}\in\{0,1,2\}$ with total $Z_c$
and an internal integer flow
\begin{equation}
 \diver F=\nu-q^{(c)},\qquad
 \mathcal F_c=\sum_{e\ {\rm in\ cell}}F_e^2.
 \label{eq:cell-flow}
\end{equation}
A spanning tree provides it; a one-site cell has zero right side.
Set $\mu=m+3r/a$ and $\rho=Z_c/V_c$.

\begin{proposition}[Fixed-cell relative bounds]
\label{prop:cell}
For $N_c=(n/p)^3$, $K=K_cN_c$ and $Z_{\rm tot}=Z_cN_c$,
\begin{equation}
 0\le E_\nu^{\rm per}(n)-E_0^{\rm per}(n)-mZ_{\rm tot}
 \le C_{\rm cell}K,
 \label{eq:cell-bound}
\end{equation}
where
\begin{equation}
 C_{\rm cell}
 =\frac{3rZ_c+(6+6r+3/g^2)V_c+(g^2/2)\mathcal F_c}{aK_c}.
 \label{eq:cell-constant}
\end{equation}
The constant may depend on the cell and its screening flow.
Every periodic source ground state also satisfies
\begin{equation}
 \langle H_E\rangle\le A_c n^3,\qquad
 A_c=2(h_*-\mu)+\mu\rho+c_E\mathcal F_c/V_c+3/(g^2a).
 \label{eq:electric-budget}
\end{equation}
\end{proposition}
\begin{proof}
Repeat the cell flow, put intercell links in electric flux zero,
and occupy the two negative-$\beta$ spinors plus
$q_x^{(c)}$ positive-$\beta$ spinors at each site.
All Gauss constraints hold and hopping expectations vanish.
Its periodic energy is at most
\begin{equation}
 -2\mu n^3+\mu Z_{\rm tot}
                 +c_EN_c\mathcal F_c+3n^3/(g^2a).
 \label{eq:cell-trial}
\end{equation}
Subtract $E_0^{\rm per}\ge-2h_*n^3$ and $mZ_{\rm tot}$
to obtain the upper bound. The lower bound is
\eqref{eq:main}. Finally $H_B\ge0$ and the fermion lower bound
give $\langle H_E\rangle\le E_\nu^{\rm per}+2h_*n^3$;
insert \eqref{eq:cell-trial}.
\end{proof}

The budget \eqref{eq:electric-budget} is absolute and extensive,
not a bound on electric excitations relative to the vacuum.
On an open rectangular union of complete cells, the same trial
gives
\begin{equation}
 E_Y^{\rm op}(\nu)\le
 \left[-2m+m\rho+c_E\mathcal F_c/V_c+3/(g^2a)\right]|Y|.
 \label{eq:open-cell-trial}
\end{equation}
Indeed the site's diagonal expectation is
$(q_x^{(c)}-2)(m+r\deg_Y(x)/(2a))\le(q_x^{(c)}-2)m$.

\begin{theorem}[Open zero-flux source limit]
\label{thm:open-source}
For each fixed cell as above,
\begin{equation}
 e_\nu^{\rm op}
 =\lim_{\substack{n\to\infty\\p\mid n}}
       \frac{E_{Q_n}^{\rm op}(\nu)}{n^3}
 =\inf_{\substack{\ell\in p\mathbb N\\\ell\ge3}}
   \left\{\frac{E_{Q_\ell}^{\rm op}(\nu)}{\ell^3}
                      +\frac{3B}{\ell}\right\}
 \label{eq:open-source-limit}
\end{equation}
exists and is finite. In addition,
\begin{equation}
 e_{\vac}+m\rho
 \le\liminf_{\substack{n\to\infty\\p\mid n}}\frac{E_\nu^{\rm per}(n)}{n^3}
 \le\limsup_{\substack{n\to\infty\\p\mid n}}\frac{E_\nu^{\rm per}(n)}{n^3}
 \le e_\nu^{\rm op}.
 \label{eq:periodic-bracket}
\end{equation}
\end{theorem}
\begin{proof}
Take products of zero-exterior-flux physical block ground states
and give newly joining rotors electric flux zero. This state
satisfies all source Gauss constraints on the joined graph.
Lemma~\ref{lem:cut} gives the upper almost-subadditivity
\[
 E_{Q_{k\ell}}^{\rm op}(\nu)
 \le k^3E_{Q_\ell}^{\rm op}(\nu)+3B(k\ell)^3/\ell.
\]
For general $n,\ell\in p\mathbb N$, the rectangular remainders
are unions of whole admissible cells. Their trial energy is
$O(\ell n^2)$ by \eqref{eq:open-cell-trial}. Consequently
the limsup of $E_{Q_n}^{\rm op}(\nu)/n^3$ is at most
$E_{Q_\ell}^{\rm op}(\nu)/\ell^3+3B/\ell$ for every fixed
compatible $\ell$. Let $\gamma$ be their infimum.
By definition $E_{Q_n}^{\rm op}(\nu)/n^3+3B/n\ge\gamma$;
hence the liminf is at least $\gamma$ and the limsup at most
$\gamma$. The extensive lower bound makes the limit finite.
Only one-sided almost-subadditivity was used.

Adjoining zero-flux closing rotors gives
$E_\nu^{\rm per}(n)\le E_{Q_n}^{\rm op}(\nu)+3Bn^2$.
This proves the last inequality in \eqref{eq:periodic-bracket}.
Its first follows from the charge-cost and vacuum-limit theorems.
\end{proof}

Thus the \emph{open-boundary} relative energy per nucleus has
limit $(V_c/K_c)(e_\nu^{\rm op}-e_{\vac})$.
A fixed subtraction $S_\infty$ subtracts its fixed density.
The periodic subtraction $S_n$ has the same limiting density:
$\G_n(0)\to\G_\infty(0)$ by ordinary Fourier Riemann convergence
away from zero and the uniform $O(\delta)$ shell bound within
$0<|k|\le\delta n$, followed by $\delta\downarrow0$.

\subsection{The periodic-source question}

We state precisely why \eqref{eq:periodic-bracket} is not yet
a periodic-source limit theorem. For a cut set $\mathcal C$,
orient each cut link once, write $f\in\mathbb Z^{\mathcal C}$,
and let $b_B(f)_x$ be its signed outgoing flux at a boundary
site $x$ of block $B$. Internal Gauss law becomes
$\diver_BE+q=\nu|_B-b_B(f)$.
Diagonalizing the retained free cut rotors gives the exact comparison
\begin{equation}
 \left|E_\nu-\inf_f
 \left\{\sum_B E_B^{\rm op}(\nu|_B-b_B(f))
                           +c_E\sum_{e\in\mathcal C}f_e^2\right\}\right|
 \le B|\mathcal C|.
 \label{eq:flux-decomposition}
\end{equation}
An empty block sector has energy $+\infty$.
At fixed $f$ the physical space of the cut Hamiltonian is the
tensor product of the displayed block sectors. The direct sum
over $f$ gives its minimum, and the bounded-cut estimate proves
\eqref{eq:flux-decomposition}. This step retains the constraints;
replacing all block energies by their vacuum minima would lose
the source comparison.

A sufficient additional hypothesis on commensurate cubes is
\begin{equation}
 E_B^{\rm op}(\nu-b_B(f))\ge E_B^{\rm op}(\nu)
      -C_\partial|\partial_{\rm links}B|
      -C_f\sum_{e\in\partial_{\rm links}B}f_e^2,
 \tag{BF}\label{eq:bf}
\end{equation}
with finite constants independent of side and flux, whenever
the left sector is admissible. Boundary links here are counted
as incidences, so a link shared by two blocks occurs twice.
Hypothesis (BF) is \emph{not proved}.

For clarity, its sufficiency can be checked without a restrictive
smallness condition on $C_f$. Average the grid of $\ell$-cubes
over translations by whole period cells. If $\ell\mid n$, a
given link is cut with probability at most $p/\ell$.
By \eqref{eq:electric-budget}, some grid therefore obeys
\[
 \left\langle\sum_{\rm cut}E_e^2\right\rangle
       \le pA_cn^3/(c_E\ell).
\]
There are $3n^3/\ell$ cut links. Applying (BF) in the state's
flux decomposition would give
\begin{equation}
 \frac{E_\nu^{\rm per}(n)}{n^3}\ge
 \frac{E_{Q_\ell}^{\rm op}(\nu)}{\ell^3}
 -\frac{3B+6C_\partial}{\ell}
 -\frac{(2C_f-c_E)_+pA_c}{c_E\ell}.
 \label{eq:conditional}
\end{equation}
For general compatible $n$, use whole-cell rectangular
remainders and average the global origin. Cut probability is
at most $p/\ell+p/n$; the remaining $O(\ell n^2)$ sites have
the absolute extensive lower bound. Thus \eqref{eq:conditional}
persists for the liminf at fixed $\ell$. Sending $\ell\to\infty$
would prove equality of the periodic and open limits.
Any finite $C_f$ suffices. Direct absorption without averaging
would instead require $2C_f\le c_E$, because of double incidence.

Admissibility in such a comparison matters. For example take
$\nu_x=4$ on even $x_1$ layers and zero on odd layers. Its mean
is the admissible value two per site. An odd $\ell$-cube starting
on an even layer has total source charge
$2\ell^3+2\ell^2>2|B|$, so its zero-flux sector is empty.
Yet $q_x=2$, internal $E_{x,1}=(-1)^{x_1}$, other components
zero, and outward flux $+1$ on both $x_1$ faces give a
finite-energy state. No finite boundary cost compares this
state to the empty zero-flux sector. Whole-cell cubes avoid
this obstruction but do not by themselves establish (BF).

The averaging estimate controls exposed electric energy; it
does not repair the block Gauss sectors. The Graf--Schenker
continuum Coulomb localization inequality, as stated in
\cite[Theorem 2]{ThermoII}, likewise does not supply (BF) for
compact integer fluxes. The periodic-source limit remains
conditional on a screening or gluing input of this kind.

\subsection{Limits of the conclusions}

The results are for a fixed spatial regulator with nonnegative
\emph{bare} Wilson mass. They do not establish that a
renormalization trajectory with fixed physical parameters stays
in that range. The self-energy bound diverges as $a^{-1}$.
There is no construction of renormalized continuum QED or
cutoff-uniform stability constant.
Compact gauge rotors provide a different finite regulator from
a transverse photon Fock-space momentum cutoff. No Coulomb-phase
or gapless-photon assertion is needed.

The electric kernel used here is that of the unshifted rotor.
Additional flux projections, nonintegral background electric
shifts, or finite electric truncations need not preserve the
positive measure in the proof and are not covered.
The charge gap compares different total-charge sectors; binding,
neutral pair gaps and crystal equilibrium require further
arguments. Finally, thermodynamic limits for the vacuum and for
open zero-flux sources do not imply the unproved periodic-source
limit. These distinctions are part of the theorem's scope.

\begin{thebibliography}{9}
\bibitem{FSS}
J. Fr\"ohlich, B. Simon, and T. Spencer,
\emph{Infrared bounds, phase transitions and continuous symmetry breaking},
Commun. Math. Phys. \textbf{50} (1976), 79--95.
\href{https://doi.org/10.1007/BF01608557}{doi:10.1007/BF01608557}.

\bibitem{Seiler}
E. Seiler, \emph{Upper bound on the color-confining potential},
Phys. Rev. D \textbf{18} (1978), 482--483.
\href{https://doi.org/10.1103/PhysRevD.18.482}{doi:10.1103/PhysRevD.18.482}.

\bibitem{KS}
J. Kogut and L. Susskind,
\emph{Hamiltonian formulation of Wilson's lattice gauge theories},
Phys. Rev. D \textbf{11} (1975), 395--408.
\href{https://doi.org/10.1103/PhysRevD.11.395}{doi:10.1103/PhysRevD.11.395}.

\bibitem{Zache}
T. V. Zache, F. Hebenstreit, F. Jendrzejewski, M. K. Oberthaler,
J. Berges, and P. Hauke,
\emph{Quantum simulation of lattice gauge theories using Wilson fermions},
Quantum Sci. Technol. \textbf{3} (2018), 034010.
\href{https://doi.org/10.1088/2058-9565/aac33b}{doi:10.1088/2058-9565/aac33b}.

\bibitem{Wei}
Z. C. Wei, C. Wu, Y. Li, S. Zhang, and T. Xiang,
\emph{Majorana Positivity and the Fermion Sign Problem of Quantum
Monte Carlo Simulations},
Phys. Rev. Lett. \textbf{116} (2016), 250601.
\href{https://doi.org/10.1103/PhysRevLett.116.250601}{doi:10.1103/PhysRevLett.116.250601}.

\bibitem{Miyao}
T. Miyao,
\emph{Upper bounds on the charge susceptibility of many-electron
systems coupled to the quantized radiation field},
Lett. Math. Phys. \textbf{105} (2015), 1119--1133.
\href{https://doi.org/10.1007/s11005-015-0775-9}{doi:10.1007/s11005-015-0775-9}.

\bibitem{ThermoI}
C. Hainzl, M. Lewin, and J. P. Solovej,
\emph{The thermodynamic limit of quantum Coulomb systems.
Part I. General theory},
Adv. Math. \textbf{221} (2009), 454--487.
\href{https://doi.org/10.1016/j.aim.2008.12.010}{doi:10.1016/j.aim.2008.12.010}.

\bibitem{ThermoII}
C. Hainzl, M. Lewin, and J. P. Solovej,
\emph{The thermodynamic limit of quantum Coulomb systems.
Part II. Applications},
Adv. Math. \textbf{221} (2009), 488--546.
\href{https://doi.org/10.1016/j.aim.2008.12.011}{doi:10.1016/j.aim.2008.12.011};
\href{https://arxiv.org/abs/0806.1709}{arXiv:0806.1709}.
\end{thebibliography}

\end{document}
